\documentclass[11pt,a4paper]{article}
\usepackage[margin=1in]{geometry}
\usepackage[T1]{fontenc}
\usepackage[utf8]{inputenc}
\usepackage{lmodern}
\usepackage{amsmath}
\usepackage{booktabs}
\usepackage{hyperref}
\usepackage{enumitem}
\usepackage{graphicx}
\usepackage{xcolor}
\setlist{itemsep=0.3em, topsep=0.3em}

\title{New Developments in the GDIS Simulation Package:\\Integration of Qbox}
\author{Author 1 \and Author 2 \and Author 3}
\date{\today}

\begin{document}
\maketitle

\begin{abstract}
Graphical workflow integration remains a practical bottleneck for first-principles studies that require rapid model setup, repeated electronic-structure runs, and continuous interpretation of intermediate results.
Building on prior GDIS external-code integration work, this paper presents a new Qbox integration path in GDIS that combines structure export, executable orchestration, pseudopotential management, asynchronous run monitoring, and restart/output re-import into a single interactive environment.
In addition to a guided setup mode, we introduce an expert pass-through mode that preserves access to advanced Qbox commands without forcing one-by-one GUI hardcoding.
The implementation supports runnable input generation, XML-based restart loading, extraction of key run observables (including total energy and maximum force), and profile presets for common workflows (SCF, geometry optimization, frozen phonons, band-structure setup, and HOMO--LUMO analysis).
Validation is shown through reproducible round-trip tests and chemistry-oriented use cases.
This integration is designed to reduce workflow friction while retaining flexibility for advanced users.
\end{abstract}

\section{Introduction}
GDIS has long served as a visualization and pre/post-processing environment for atomistic simulation.
Recent development cycles have focused on renewing code paths, modernizing rendering/UI behavior, and extending practical interoperability with external engines.
While previous GDIS integration work emphasized VASP and USPEX, Qbox remained underrepresented in integrated GUI-to-run workflows.

Qbox itself exposes a rich command set and variable space for electronic-structure optimization, ionic dynamics, response workflows, and post-processing.
However, many users still rely on fragmented command-line scripts for run preparation and result inspection, especially in teaching and exploratory workflows.
The goal of this work is to provide a Qbox path in GDIS that is simple for common tasks, while preserving access to advanced command-level control.

\section{Design Goals}
The implementation was guided by four requirements:
\begin{enumerate}
  \item \textbf{Round-trip operability}: create runnable Qbox inputs from active GDIS models, run Qbox, and reload saved outputs in one loop.
  \item \textbf{Practical defaults}: provide robust default settings and species-path handling for rapid startup.
  \item \textbf{Expert extensibility}: avoid GUI lock-in by allowing arbitrary Qbox command insertion and include-file workflows.
  \item \textbf{Traceable execution}: preserve log files and asynchronous task visibility through existing GDIS task infrastructure.
\end{enumerate}

\section{Implementation}
\subsection{Qbox File I/O Layer}
We implemented a dedicated Qbox I/O layer with:
\begin{itemize}
  \item structure-to-input writer support (\texttt{cell}, \texttt{species}, \texttt{atom});
  \item direct parser support for Qbox-like input files;
  \item restart/output XML ingestion for structure updates and key observables.
\end{itemize}

For non-periodic models, export logic automatically places the structure in a vacuum box.
For restart/output parsing, atom positions are updated and scalar observables (e.g., final total energy and max force) are mapped into model properties for inspection.

\subsection{Qbox Runtime Integration in GDIS}
The Qbox dialog integrates model metadata, working-directory paths, default run controls, and pseudopotential file assignment.
Runtime input generation enforces a runnable script path and validates pseudopotential URI existence before run submission.
When valid, jobs are launched asynchronously through the existing task manager and write run logs to user-specified locations.

\subsection{Expert Pass-Through Mode}
A major extension in this work is an expert mode that supports command-level composition without waiting for dedicated GUI widgets for each Qbox feature.
The expert mode includes:
\begin{itemize}
  \item optional model export block enable/disable;
  \item optional default settings block enable/disable;
  \item pre-default and post-default free-form command blocks;
  \item optional include-command-file injection (e.g., displacement scripts such as \texttt{moves.i});
  \item auto-save and auto-quit toggles with duplicate-command suppression.
\end{itemize}

This design allows users to directly pass advanced commands (e.g., \texttt{kpoint}, \texttt{set scf\_tol}, \texttt{set nempty}, \texttt{compute\_mlwf}, \texttt{response}, \texttt{spectrum}, custom \texttt{run} schedules) while keeping an accessible GUI entry point.

\subsection{Preset Profiles}
To bridge novice and advanced usage, preset profiles auto-fill expert fields for common workflows:
\begin{itemize}
  \item SCF
  \item Band
  \item GeoOpt
  \item FrozenPhonon
  \item HOMO--LUMO
\end{itemize}
Presets populate fields but do not auto-run; users review and edit before execution.

\section{Validation and Demonstration Plan}
\subsection{Functional Validation}
Minimum reproducible checks include:
\begin{enumerate}
  \item build validation for GTK4 primary and GTK2 fallback targets;
  \item Qbox round-trip smoke test from generated input to saved XML re-import;
  \item optional dependency detection and executable-path verification.
\end{enumerate}

\subsection{Scientific Workflow Demonstrations}
Suggested demonstration cases:
\begin{itemize}
  \item methane SCF convergence (\texttt{ecut} sweep and \texttt{etotal} extraction);
  \item geometry optimization with force threshold tracking;
  \item frozen-phonon displacement scripting via include-file mode;
  \item HOMO--LUMO setup and wavefunction-plot command generation;
  \item monolayer band-structure setup using k-point command injection.
\end{itemize}

\subsection{Proposed Figures and Tables}
\textbf{Figure candidates}
\begin{itemize}
  \item Qbox dialog: Setup, Potentials, and Expert tabs;
  \item round-trip workflow diagram (GDIS model $\rightarrow$ Qbox run $\rightarrow$ XML reload);
  \item representative model reloaded from Qbox XML;
  \item log/status view during asynchronous execution.
\end{itemize}

\textbf{Table candidates}
\begin{itemize}
  \item capability matrix: guided mode vs expert mode coverage;
  \item validation checklist and outcomes;
  \item mapping of paper use cases to preset profiles.
\end{itemize}

\section{Discussion}
The integration prioritizes a hybrid architecture: guided defaults for speed and expert pass-through for completeness.
This approach scales better than hardcoding every Qbox command into separate widgets and reduces maintenance overhead as upstream Qbox evolves.

Current limitations include incomplete command-semantic validation in expert free-form blocks and remaining UI modernization tasks outside the Qbox workflow.
Future work should add stronger command linting, richer result parsing (band/eigenvalue extraction helpers), and automated regression tests for command templates.

\section{Conclusion}
We introduced a practical and extensible Qbox integration path in GDIS that supports end-to-end run orchestration, output re-ingestion, and expert-level command flexibility.
By combining structured defaults with pass-through control and reusable presets, the implementation improves usability without reducing access to advanced Qbox workflows.
This work provides a concrete foundation for both teaching and research pipelines that rely on Qbox within a graphical simulation environment.

\section*{Reproducibility Notes}
\begin{itemize}
  \item Record repository URL and commit hash in final camera-ready manuscript.
  \item Archive benchmark inputs, include files, and generated logs used for figures/tables.
  \item Include exact build target (GTK4 or GTK2) in Methods.
\end{itemize}

\section*{References}
\begin{enumerate}
  \item H. Okadome Valencia, B. Wang, G. Frapper, A. L. Rohl, \emph{New developments in the GDIS simulation package: Integration of VASP and USPEX}, J. Comput. Chem. 42 (2021) 1602--1626. doi:10.1002/jcc.26697.
  \item Qbox Documentation, \url{http://www.qboxcode.org/doc/html/index.html}.
  \item Qbox Command Reference, \url{http://qboxcode.org/doc/html/usage/commands.html}.
  \item Qbox Variables Reference, \url{http://qboxcode.org/doc/html/usage/variables.html}.
\end{enumerate}

\end{document}
